\documentclass[11pt]{article}

\usepackage[margin=1in]{geometry}
\usepackage{amsmath, amssymb, amsfonts}
\usepackage{booktabs}
\usepackage{enumitem}
\usepackage{xcolor}
\usepackage{listings}
\usepackage[colorlinks=true, linkcolor=blue!60!black, urlcolor=blue!60!black, citecolor=blue!60!black]{hyperref}
\usepackage{parskip}

\lstset{
  basicstyle=\ttfamily\small,
  breaklines=true,
  frame=single,
  framesep=4pt,
  backgroundcolor=\color{gray!8},
  showstringspaces=false
}

\newcommand{\x}{\mathbf{x}}
\newcommand{\vv}{\mathbf{v}}
\newcommand{\g}{\mathbf{g}}
\newcommand{\R}{\mathbf{R}}
\newcommand{\tr}{\mathbf{t}}

\title{\textbf{PhysSplat: A Physically Grounded 3D Neural World Model \\ for Interactive Environments}\\[6pt]
\large Technical Design Document}
\author{Mohamad Salman}
\date{August 2026}

\begin{document}
\maketitle

\begin{abstract}
PhysSplat converts a photograph of everyday rigid objects on a flat surface (blocks, pencils, erasers, and similar simple items) into an interactive, physically plausible 3D simulation that runs in a web browser. The pipeline lifts the image into a 3D Gaussian Splat representation, decomposes the splat cloud into discrete rigid bodies without supervision, and simulates those bodies with a learned graph-network dynamics model trained entirely on synthetic data. The user can orbit the scene, poke and shove objects with the mouse, and watch them topple, slide, and roll, with every frame rendered from the same Gaussians that were reconstructed from the photo. This document specifies the full architecture, the mathematics of the learned simulator, the synthetic data pipeline, and a step-by-step build plan starting from an empty repository on an Apple M4 Pro with a compute budget under \$30.
\end{abstract}

\tableofcontents
\newpage

%=====================================================================
\section{What the Project Is and What It Demonstrates}
%=====================================================================

\subsection{The one-sentence version}
Take a photo of a few simple objects on a table (blocks, pencils, erasers), get back a browser window where those exact objects exist in 3D and obey physics you can interact with, where the physics is a neural network rather than a hand-coded solver. The scope is deliberately bounded to \emph{convex-ish rigid objects on a flat surface}: things that stack, slide, tip, and roll. Concave or articulated objects (mugs, scissors) are out of scope.

\subsection{Why this is interesting}
Two mature technologies exist on either side of this project, and neither crosses the gap alone:

\begin{itemize}[itemsep=2pt]
  \item \textbf{Generative 3D reconstruction} (3D Gaussian Splatting, feed-forward image-to-3D models) produces photorealistic 3D scenes from images, but the output is a hollow visual shell. Nothing has mass, nothing collides, gravity does not exist.
  \item \textbf{Physics engines} (MuJoCo, PyBullet, Rapier) simulate rigid bodies robustly, but they require clean meshes, convex decompositions, and hand-tuned parameters. They cannot ingest a photograph.
\end{itemize}

PhysSplat bridges the two with a \emph{learned} dynamics model: a graph neural network trained on synthetic trajectories that predicts how bodies move under gravity, contact, friction, and user-applied impulses. No meshing, convex decomposition, or manual parameter tuning is needed at inference time.

\subsection{The world model formulation}
The system learns an implicit transition function
\begin{equation}
  s_{t+1} = f_\theta(s_t, a_t),
\end{equation}
where the state $s_t$ is a set of particles with positions, velocities, and body assignments; the action $a_t$ is an external impulse injected by the user at a 3D location; and $f_\theta$ is a graph neural network that models contact through message passing rather than analytic collision resolution. This is the defining property of a world model: the dynamics are learned from data, not programmed.

\subsection{What it demonstrates to a technical reviewer}
\begin{itemize}[itemsep=2pt]
  \item \textbf{Spatial intelligence:} lifting unstructured 2D pixels into a structured, object-decomposed 3D representation.
  \item \textbf{Geometric deep learning:} a graph network simulator with relational inductive biases, rigid-body structure, noise-robust multi-step rollout, and correct input/output normalization.
  \item \textbf{ML systems engineering:} synthetic data generation at scale, training on Apple Silicon via Metal Performance Shaders, ONNX export, and client-side neural network inference in the browser at interactive rates.
  \item \textbf{Full-stack execution:} a deployed Next.js application with a real-time WebGL Gaussian renderer.
\end{itemize}

%=====================================================================
\section{System Architecture}
\label{sec:architecture}
%=====================================================================

\begin{center}
\begin{tabular}{@{}lll@{}}
\toprule
\textbf{Stage} & \textbf{Runs} & \textbf{Frequency} \\
\midrule
A. Image $\rightarrow$ 3D Gaussians & Server / local Mac & once per scene \\
B. Decomposition into rigid bodies & Server / local Mac & once per scene \\
C. Physics particle resampling \& canonicalization & Server / local Mac & once per scene \\
D. Learned dynamics rollout (GNN) & \textbf{Browser (ONNX Runtime Web)} & 60\,Hz \\
E. Gaussian rendering \& interaction & Browser (WebGL/WebGPU) & 60\,Hz \\
\bottomrule
\end{tabular}
\end{center}

Data flow: a photo is uploaded; Stage A produces a splat file; Stage B labels every Gaussian with a body id and extracts the ground plane; Stage C samples physics particles per body and normalizes the scene (gravity-aligned, metric-scaled); the browser receives one packet containing the labeled Gaussians, the physics particles, and body metadata, and from then on everything runs client-side: the GNN steps the particles, rigid transforms are recovered per body, the Gaussians are moved, sorted, and splatted.

Two design decisions shape everything downstream:

\textbf{Two coupled representations.} The physics does not run on the raw Gaussians. A reconstruction produces tens of thousands of nonuniformly dense Gaussians, which is too slow for a radius-graph GNN and statistically unlike any training data. Instead, each body gets a small set of uniformly resampled \emph{physics particles} (Section~\ref{sec:stageC}) that the GNN simulates, and the full Gaussian set is carried along rigidly: each frame, body $b$'s transform $(\R_b, \tr_b)$ is applied to its Gaussians as $\boldsymbol{\mu}' = \R_b \boldsymbol{\mu} + \tr_b$ and $\boldsymbol{\Sigma}' = \R_b \boldsymbol{\Sigma} \R_b^\top$.

\textbf{Client-side inference.} Round-tripping physics steps through a server adds 50 to 300\,ms of latency per step and destroys interactivity. The dynamics network is small (about 5M parameters over about 1{,}000 particles), so it runs in the browser via ONNX Runtime Web at a few milliseconds per step. The server's only job is the one-time reconstruction, which is latency-tolerant.

Training happens entirely offline (Section~\ref{sec:training}) on synthetic PyBullet scenes that mimic Stage C's output format, which is what makes the sim-to-recon transfer work.

%=====================================================================
\section{The Pipeline, Stage by Stage}
%=====================================================================

\subsection{Stage A: Image to 3D Gaussians}
\label{sec:stageA}

\textbf{Primary path (single image).} A feed-forward image-to-3D model runs once per scene:
\begin{itemize}[itemsep=2pt]
  \item \textbf{LGM}~\cite{lgm}: generates multi-view images with a diffusion model, then regresses Gaussians directly. Native Gaussian output. Heavier, and the multi-view diffusion step is the flakiest part on MPS; if MPS misbehaves, run this one stage on a rented GPU or free Colab, since it is offline and one-shot.
  \item \textbf{TripoSR}~\cite{triposr}: faster and more robust, but outputs a triplane field, so we extract a mesh (marching cubes), sample its surface, and instantiate one small isotropic Gaussian per sample with the surface color. Slightly ``point-cloudy'' visuals but serviceable.
\end{itemize}
Both are trained on single centered objects. A compact tabletop arrangement of a few small objects against a plain background reads as one object, which these models handle acceptably; the occluded back is hallucinated, but simple convex shapes make this rarely noticeable. Background removal (\texttt{rembg}) is applied before reconstruction. One caution: very thin structures (a lone pencil) are the weakest case for feed-forward reconstructors and may come back lumpy; the physics tolerates this (the convex hull of a lumpy pencil is still pencil-shaped), but for hero shots prefer chunkier objects or the multi-view path.

\textbf{Fallback path (multi-view).} 30 to 60 phone photos orbiting the scene, processed by a standard 3DGS pipeline~\cite{3dgs}, yield a far higher-fidelity splat scene including a real ground plane. The downstream pipeline is identical.

\textbf{Output:} a \texttt{.ply} of $N$ Gaussians, each with position $\boldsymbol{\mu}_i \in \mathbb{R}^3$, covariance (stored as scale and quaternion), opacity, and spherical-harmonic color coefficients.

\subsection{Stage B: Unsupervised decomposition into bodies}
\label{sec:stageB}

\begin{enumerate}[itemsep=2pt]
  \item \textbf{Gravity alignment and scaling.} Align the dominant plane normal to $+z$ and rescale so object sizes are plausible in metric units (the object class is known, so metric scale is assumed rather than estimated).
  \item \textbf{Ground extraction.} RANSAC plane fit on the lowest 30\% of points (by $z$). Inliers are labeled \emph{static ground}. If the single-image path reconstructs no ground, a synthetic plane is added at the scene's minimum $z$.
  \item \textbf{Clustering.} HDBSCAN over $[\,\boldsymbol{\mu}_i / \sigma_{xyz},\; \lambda\,\mathbf{c}_i\,]$, normalized position concatenated with weighted RGB from the DC spherical-harmonic term. Position separates spatially distinct objects; color separates touching objects of different colors. Clusters below a minimum size are merged or discarded as floaters. Because real objects are multi-colored (a pencil has a colored barrel, a lighter grip, an eraser, a tip), an \emph{over-segmentation merge pass} follows: adjacent clusters whose principal axes are collinear are combined into one body.
  \item \textbf{Per-body statistics.} Centroid, principal axes (PCA), an oriented bounding box (used for mouse-ray picking), and mass plus inertia estimated from the cluster's convex hull with an assumed density.
\end{enumerate}

\textbf{Output:} a body id per Gaussian, a ground label, and per-body metadata.

\subsection{Stage C: Physics particle resampling and canonicalization}
\label{sec:stageC}

This stage makes real reconstructions statistically indistinguishable from training data, which is the entire trick behind sim-to-real transfer here.

For each body, sample surface points by farthest-point sampling at a \emph{fixed target spacing} of $\sim$4\,mm (capped at $P_{\max} \approx 400$ per body), rather than a fixed count. Fixed spacing keeps neighborhood statistics consistent across object sizes: a pencil and a large block have the same local particle density even though their counts differ, so the contact radius stays a single global constant. The sampler must be \emph{byte-for-byte the same procedure} used when generating training data. Each particle stores its offset from the body centroid in the body frame, which is what lets us recover the body's rigid transform. Ground is handled analytically as a node feature rather than with ground particles.

\textbf{Output:} particle positions $\x_i$, body ids $b_i$, body-frame offsets $\mathbf{r}_i$, per-body mass and inertia, and zero initial velocities.

%=====================================================================
\section{The Learned Dynamics Engine}
\label{sec:dynamics}
%=====================================================================

The core of the project: an Encode--Process--Decode graph network in the GNS family~\cite{gns}, with a rigid-body output head. Classic GNS predicts an independent acceleration per particle, which works for sand and water but fails for rigid bodies: accumulated per-particle error violates rigidity within tens of steps and objects visibly melt. PhysSplat therefore keeps the particle-level contact graph but constrains the output to rigid motion, following the per-body approach of FIGNet~\cite{fignet}.

\subsection{Graph construction (every step)}
Nodes are the physics particles. Edges connect pairs within a contact radius:
\begin{equation}
\mathcal{E}_t = \{(i,j) : \lVert \x_i - \x_j \rVert < r_{\text{contact}},\; i \neq j\},
\end{equation}
with $r_{\text{contact}}$ about 1.5$\times$ the particle spacing. Same-body edges propagate information across a body; cross-body edges carry contact, and an edge feature flags which kind each edge is. Neighbor search uses a uniform grid hash, $O(N)$ per step. Edges are formed from the union of current and velocity-extrapolated positions ($\x + \vv\,\Delta t$), so a fast approach is sensed one step early instead of crossing the sensing shell unnoticed between steps; training-time graphs are built with the identical rule so there is no train/inference mismatch.

\subsection{Input features and normalization}
Per-node input:
\begin{equation}
\mathbf{h}_i^{(0)} = \big[\; \vv_i^{t-C+1}, \ldots, \vv_i^{t} \;\big\|\; d_i^{\text{ground}} \;\big\|\; m_{b_i} \;\big\|\; \operatorname{diag}(\mathbf{I}_{b_i}) \;\big\|\; \mathbf{a}_i^{\text{ext}} \;\big],
\end{equation}
a history of $C=5$ finite-difference velocities, a clipped distance to the ground plane $d_i^{\text{ground}} = \min(z_i, r_{\text{contact}})$ (how the network senses the floor without floor particles), the body mass, the diagonal of the body's inertia tensor in its principal frame (a pencil and a cube respond very differently to the same push; the network could infer this from geometry, but providing it is cheap and helps), and the external impulse feature $\mathbf{a}_i^{\text{ext}}$ (zero except during user interaction, Section~\ref{sec:interaction}). Absolute positions are never input, giving translation invariance for free.

Per-edge input for $(i,j)$: relative displacement $\x_j - \x_i$, its norm, and the same-body flag.

Every input and target dimension is standardized by dataset statistics. In GNS-class models this is not a nicety; it is the difference between converging and not, because raw accelerations are dominated by gravity while contact spikes are orders of magnitude rarer.

\subsection{Encoder, processor, decoder}
\textbf{Encoder:} two-layer MLPs lift node and edge features to $d=128$-dimensional latents.

\textbf{Processor:} $L=10$ message-passing blocks with residual connections:
\begin{align}
\mathbf{e}_{ij}^{(\ell+1)} &= \mathbf{e}_{ij}^{(\ell)} + \phi_e\!\big(\mathbf{e}_{ij}^{(\ell)}, \mathbf{h}_i^{(\ell)}, \mathbf{h}_j^{(\ell)}\big), \\
\mathbf{h}_i^{(\ell+1)} &= \mathbf{h}_i^{(\ell)} + \phi_v\!\Big(\mathbf{h}_i^{(\ell)}, \textstyle\sum_{j \in \mathcal{N}(i)} \mathbf{e}_{ij}^{(\ell+1)}\Big),
\end{align}
where $\phi_e, \phi_v$ are two-layer MLPs with LayerNorm. Ten hops let a contact on one corner of a body influence the opposite corner within a single step. Implementation is plain PyTorch (edge MLPs on gathered node pairs, aggregation via \texttt{index\_add\_}); PyTorch Geometric is deliberately avoided because its scatter kernels have poor MPS coverage and complicate ONNX export.

\textbf{Decoder (per-body rigid head):} mean-pool final node latents within each body $b$ and decode
\begin{equation}
(\Delta \vv_b,\; \Delta \boldsymbol{\omega}_b) = \psi\Big( \tfrac{1}{|b|} \textstyle\sum_{i \in b} \mathbf{h}_i^{(L)} \Big) \in \mathbb{R}^6,
\end{equation}
a linear and angular acceleration at the body's center of mass. Because every particle of a body moves under one SE(3) update, rigidity is exact by construction and rollouts cannot melt.

\subsection{Integration}
Semi-implicit Euler per body. Everything analytically knowable is integrated exactly: gravity, and the user-applied force $\mathbf{F}_b^{\text{ext}}$ whose vector and application point $\mathbf{p}_b$ are both known, so its full Newton--Euler effect on a free body is computed in closed form. The network learns only the residual that cannot be written down: contact and friction.
\begin{align}
\vv_b^{t+1} &= \vv_b^{t} + \Big(\g + \tfrac{\mathbf{F}_b^{\text{ext}}}{m_b}\Big) \Delta t + \Delta \vv_b, &
\boldsymbol{\omega}_b^{t+1} &= \boldsymbol{\omega}_b^{t} + \mathbf{I}_b^{-1}\big((\mathbf{p}_b - \bar{\x}_b) \times \mathbf{F}_b^{\text{ext}}\big) \Delta t + \Delta \boldsymbol{\omega}_b, \\
\tr_b^{t+1} &= \tr_b^{t} + \vv_b^{t+1} \Delta t, &
\R_b^{t+1} &= \exp\!\big([\boldsymbol{\omega}_b^{t+1}]_\times \Delta t\big)\, \R_b^{t},
\end{align}
and particle positions follow: $\x_i^{t+1} = \R_b^{t+1} \mathbf{r}_i + \tr_b^{t+1}$. Training targets subtract the same analytic terms, so the supervision is the pure contact residual. This construction \emph{guarantees} point-of-application physics rather than hoping the network learns it: grab a pencil off-center and it leans until its center of mass hangs under the cursor; grab an end and it dangles as a pendulum, damped by the spring's $k_d$; release it and free fall is exact. What remains learned is exactly what should be: how bodies respond to each other.

\subsection{Non-penetration in a constraint-free model}
A learned simulator has no hard collision constraint, so three stacked mechanisms keep bodies from passing through each other. First, the data and the noise: ground truth contains essentially no interpenetration, and training-noise injection deliberately produces slightly overlapping input states whose clean targets push apart, teaching active penetration \emph{recovery} rather than mere avoidance. Second, a speed budget: contact is sensed only within $r_{\text{contact}}$, so relative motion faster than roughly $r_{\text{contact}}/\Delta t \approx 0.36$\,m/s can cross the sensing shell between steps; the predictive edges above extend this budget, data generation runs PyBullet at 240\,Hz internally so ground truth itself never tunnels, and at inference user force and grab-speed caps keep relative speeds inside the budget (0.36\,m/s is already fast across a 15\,cm scene). Third, an analytic guard: every demo body is a capsule, and capsule-capsule overlap is a closed-form segment-segment distance test, so after each learned step a position-based projection resolves residual overlaps among the $\sim$15 body pairs and against the ground plane at negligible cost. The guard exists to clean up millimetres, not to do the physics: evaluation records \emph{pre-guard} penetration, and a model that leans on the guard fails the scorecard. The speed budget is nonetheless the weakest of the three, because a user is free to violate it: dropping a pencil from $250$\,mm or yanking one out from under a pile drives relative speeds well past $0.36$\,m/s, and both tunnelled by $7.8$\,mm before a swept version of the guard was added (Section~\ref{sec:invariants}).

\subsection{The shape-matching fallback}
If the per-body head underfits (pooling discards where on the body contact happened), a fallback decoder predicts per-particle accelerations as in classic GNS, integrates particles freely, then snaps each body back to rigidity via the Kabsch projection: with predicted positions $\tilde{\x}_i$ and offsets $\mathbf{r}_i$,
\begin{equation}
\R_b^\ast = \arg\max_{\R \in SO(3)} \operatorname{tr}\!\big(\R^\top \mathbf{A}_b\big), \qquad
\mathbf{A}_b = \sum_{i \in b} (\tilde{\x}_i - \bar{\x}_b)\, \mathbf{r}_i^\top,
\end{equation}
solved by SVD of $\mathbf{A}_b$, with $\tr_b^\ast = \bar{\x}_b$. Both heads share the same encoder and processor, so switching is a one-line config change.

\subsection{Loss and rollout stability}
\label{sec:loss}
Training minimizes single-step MSE on normalized accelerations. Single-step training alone drifts over long rollouts because the model never sees its own errors; two standard remedies are used:
\begin{itemize}[itemsep=2pt]
  \item \textbf{Training-noise injection:} corrupt input positions and velocities with random-walk Gaussian noise ($\sigma \approx$ 1 to 3\,mm over the history window) while computing targets against the clean next state. The model learns to \emph{correct} perturbations, exactly the skill rollout requires. This is the highest-leverage trick in the GNS literature.
  \item \textbf{Short fine-tuning rollouts:} after single-step convergence, fine-tune on 4 to 8-step unrolled losses at a small learning rate.
\end{itemize}

%=====================================================================
\section{Browser Runtime, Rendering, and Interaction}
\label{sec:stageE}
%=====================================================================

\subsection{Client-side inference}
The trained model is exported to ONNX with symbolic shapes for dynamic node and edge counts. Scatter aggregation was to export as ONNX \texttt{ScatterND}; in practice ONNX Runtime's \texttt{ScatterND} races under multiple threads and gave a $2.7\times10^{-2}$ parity error, so the export uses a race-free segmented sum instead: edges are sorted by receiver and carried with segment pointers, and the aggregation is a deterministic prefix sum ($2.2\times10^{-6}$ parity). Graph construction (grid-hash neighbour search) lives in JavaScript and rebuilds the edge list each step.

Runtime: ONNX Runtime Web proved unusable as the fast path. Its WebGPU provider dispatches $\sim$200 operators with about a millisecond of overhead each (measured 190--490\,ms per step) and cannot graph-capture this model. The shipped backend is therefore a hand-written WebGPU one: five WGSL kernels (a register-tiled GEMM, layer norm, edge gather, segmented sum, body pooling) recorded into a single compute pass per step, with cached bind groups and persistent input buffers. On an M-series laptop a four- to five-pencil scene steps in 12--45\,ms, and the ONNX Runtime WASM path remains as a correct fallback. Rendering interpolates between physics states so the display stays at 60\,Hz regardless.

\subsection{Rendering moving bodies}
The plan was to render the reconstruction directly: a Three.js scene using a 3DGS splatting library, each body's Gaussians forming their own batch with a per-batch model matrix, so the CPU uploads six floats per body per frame rather than rewriting buffers, covariances rotating by conjugation with $\R_b$, and depth sorting retriggered when bodies move. That path exists and is reachable with \texttt{?points=1}, and the reconstruction's own triangles are reachable with \texttt{?mesh=capsule}. Neither is the default, for a reason worth stating plainly: a reconstruction good enough to simulate is not good enough to look at. The splats read as point-cloudy and the extracted triangles are lumpy, so the pencils looked like approximations of pencils, and a viewer who cannot believe the object will not believe its motion either.

What ships instead is a procedural model of the object, driven by the reconstruction rather than replaced by it. A BIC Matic Grip is built from its real parts --- eraser, cap, constant-diameter barrel, clip, moulded grip, plastic cone, lead --- at proportions measured from a photograph of the pencil, since those proportions are identical across every pencil in the set. Only what actually varies is read from the reconstruction: the body's length and radius profile, and its barrel colour, taken as the most saturated cluster of the reconstructed colours rather than their mean, so the grey grip and the shaded side do not muddy it. The two are tied together deliberately. The grip is drawn at the fattest band of the particle hull the physics moves, and the barrel follows from the pencil's real $11$ over $9$\,mm ratio, so no drawn surface lies outside the simulated body. This matters more than it sounds: an earlier version drew the barrel at the capsule's \emph{median} particle radius while the physics reached the $95$th percentile, and every pencil the ground rule was holding exactly on the table was drawn $1.7$\,mm above it. Users reported that as hovering, and the same mismatch from the other side --- a clip standing $0.9$\,mm proud of the drawn barrel --- was reported as sinking into the table. The renderer and the simulator must agree on where the object's surface is, whatever draws it.

\subsection{Interaction}
\label{sec:interaction}
Mouse down casts a ray tested against per-body OBBs from Stage B. Two modes share one mechanism. A quick drag-release is a \emph{poke}: a short force burst in the drag direction. Press-and-hold is a \emph{grab}: a spring-damper attaches between the grab point and the mouse target, $\mathbf{f} = k_p(\x_{\text{target}} - \x_{\text{grab}}) - k_d\,\vv_{\text{grab}}$, applied every step until release, with the same gains used in data generation. The force's rigid-body effect is integrated analytically at the true application point (Section on integration), so grab physics is exact: click the tip and it pivots from the tip, grab off-center and it leans, hold an end and it dangles, release and it falls and bounces with plastic restitution. The force also enters the network as the per-particle channel $\mathbf{a}_i^{\text{ext}}$ with Gaussian falloff from the application point, so the contact model can anticipate its consequences (pressing a pencil into the pile). Forces and target speeds are capped to the training distribution's maxima, which also keeps relative speeds inside the tunneling budget. Sliders for gravity scale and impulse strength, and a reset button (free, since the initial state is just the Stage C packet).

\subsection{Analytic invariants the residual may not violate}
\label{sec:invariants}
A learned simulator asked to run far outside its training distribution --- photo-reconstructed pencils rather than PyBullet capsules --- asserts things that are not physics. Three invariants are therefore stated explicitly, two of them active in the default runtime and the third measured and rejected, and every intervention is recorded per step so a reader can always separate the model's answer from the correction.

\emph{Free flight.} A body with no contact edge to another body and no particle within the contact radius of the floor feels only gravity and the applied force, so its residual is set to zero. Without this rule a pencil released in mid-air hung there: the network never saw a motionless unsupported body in training, and answered ``$+g$, hold still''. Measured free-fall acceleration after the rule: $1.00\,g$.

\emph{No free energy, and why it is off.} Contact forces from static things cannot add mechanical energy to a scene, and an applied force can add only the work it does. This is true, and it is nonetheless not in the default runtime, which is worth recording because it is the clearest case where measurement overruled a principle. The rule trial-integrates each step and, if the residual would raise $\sum_b \tfrac{1}{2}m_b\|\vv_b\|^2 + \tfrac{1}{2}\boldsymbol\omega_b^{\!\top}\mathbf{I}_b\boldsymbol\omega_b + m_b g z_b$ by more than $\mathbf{f}\cdot\vv_{\text{grab}}\,\Delta t$ (plus $0.5\,\mu$J of slack for pushing out of overlaps), scales the residual down by bisection until it does not. It removes the failure it was written for: a pencil at rest rotated up to $89^\circ$ on its own and balanced on its end. But scored on the synthetic scorecard it costs more than it buys, $42.4$ composite whole-scene and $48.1$ restricted to quiescent bodies against $62.3$ for free flight alone, because a bisection that cannot tell a spurious gain from a real one also damps genuine collision response. The spontaneous rearing it was meant to stop is instead prevented by the support and pivot rule below, which is a statement about equilibrium rather than about energy and does not touch the residual elsewhere. The energy rule survives behind \texttt{--energy} so the number can be reproduced.

\emph{Support and pivot.} Support is counted only from contacts below the body (contact normal pointing up) and from the floor. A body whose centre of mass projects outside the convex hull of its support is not in equilibrium, so it is integrated as a pendulum about the nearest point of that hull's boundary, $\alpha = \|\mathbf{r}\times m\mathbf{g}\|/I_{\text{hinge}}$, until it lies flat. Without this rule a pencil that landed on one end stayed tilted with the other end in the air.

Smaller cleanups sit alongside them. A supported body that has been slow for fifteen steps is held exactly still, since the model's contact response rings for about a second after every landing. Residual capsule and floor overlap is removed geometrically, and two details of that projection turned out to matter more than the projection itself. It iterates until the worst overlap in the scene falls below $0.3$\,mm rather than running a fixed number of passes, because each pass is capped at $1.5$\,mm of correction to avoid teleporting bodies, and a body driven $10$\,mm into the table therefore needs seven passes; the loop's stopping test takes the worst of body-body and body-ground overlap, and an earlier version that measured only the former would exit after one pass with a pencil still under the table. Alongside it, a swept test samples the segment each body travelled during the step and rejects any pair that deepens more than $0.3$\,mm beyond where that pair started, which is what allows it to run while a pencil being pulled from a pile is legitimately in contact with everything around it; a first version skipped pairs already touching and so disabled itself in exactly that case. Finally a held body's speed is clamped to $0.25$\,m/s after integration. Together these take a $250$\,mm drop from $7.8$\,mm of pass-through to $0.0$--$0.3$\,mm and a fast pull from $7.7$\,mm to $0.4$--$2.1$\,mm. All of it is demo-only by default; \texttt{scripts/evaluate.py} exposes the physical rules as flags so their effect on the scorecard is measured rather than assumed.

\subsection{Postscript: the demo ships a classical solver by default}
\label{sec:solver}
Everything above describes the learned simulator, and it remains available in the demo behind \texttt{?engine=gnn}. The default engine is now Rapier, a rigid-body solver compiled to WebAssembly, and it is worth stating plainly how that came about, because it is the honest result of the measurement discipline this document insists on. Eight blind review rounds and three rounds of player reports were spent adding analytic rules around the learned model's failure modes; each was scored and each earned its place; and pencils still did not lie flat, still hovered, still twitched. The change that finally made them the same object, the canonical profile of Section~\ref{sec:invariants}, is also what made a classical solver usable, since there was at last a clean collider to give it. Held to the same sensor-driven checks, the solver puts a lone dropped pencil at $0.5^\circ$ on its grip where the learned model pinned it wherever it stopped, holds a pencil still at $0.2$ to $3.8^\circ$ of tilt where the learned model reared $3$ to $35^\circ$, and costs $0.5$\,ms a step against $15$ to $33$. The learned model is kept because the comparison is the finding: a GNS-style simulator trained on synthetic capsules transfers to photographed pencils well enough to be interactive and not well enough to be right, and every analytic rule in Section~\ref{sec:invariants} is a measurement of exactly where.

\subsection{Diagnostics}
\label{sec:diagnostics}
The demo carries a diagnostic layer that records, every step: per-body pose and velocities, lowest surface point, axis elevation, capsule contacts with signed gaps, which bodies support which, penetration depth, resting state and drift since rest, kinetic energy, the model's predicted acceleration, and each rule's intervention. Twelve detectors (explosion, sinking, floating, tip balance, tilt hold, unbalanced rest, rearing, creep, rotational creep, penetration, jitter, spontaneous motion) open and close logged episodes with peaks. On top of that sit scripted probes --- rest, lift, grab at an end, poke, flick, pull the load-bearing body, drop --- which drive the same code path as the mouse and return measurements: cursor tracking error, fitted free-fall acceleration, contact onset, impact speed against the analytic value, bounces, time to rest, and neighbour displacement. The point is that a tester, human or automated, reads what the pencils did instead of inferring it from pixels; every defect listed in Section~\ref{sec:invariants} was found and confirmed this way.

%=====================================================================
\section{Synthetic Training Data}
\label{sec:data}
%=====================================================================

All training data is generated locally with PyBullet on CPU, deterministic seeds. Nothing is downloaded or labeled.

\subsection{Scene distribution}
Each trajectory:
\begin{itemize}[itemsep=2pt]
  \item 2 to 6 objects drawn from a randomized \emph{primitive-shape family}: boxes (erasers, blocks), cylinders and capsules (pencils, markers, batteries), and random convex polyhedra. Longest dimension 2 to 16\,cm, aspect ratios up to 15:1 to cover pencil-like shapes; friction $\mu \in [0.2, 0.9]$ (smooth plastic barrels sit near the low end, rubber grips near the high end); restitution $\in [0.1, 0.5]$ (hard plastic on a desk bounces visibly; excluding dead-thud materials keeps the learned bounce crisp); density randomized in $[400, 900]$\,kg/m$^3$ (hollow plastic pencils have effective density near 530; mass and inertia are input features, so the network absorbs the range). Convexity is the scope boundary: any real object is simulated as its convex hull, which is accurate for the target class. Elongated cylinders use capsule collision primitives in PyBullet, since cylinder-cylinder point contacts (crossed pencils) are numerically fragile.
  \item Initial conditions mixed across six regimes: (a) stable-ish stacks (settled 100 steps before recording), (b) objects lying scattered flat on the surface, (c) random drops from low height, and four pencil-stack regimes that mirror the demo scenes: (d) parallel bundles lying flat and touching, (e) crosshatch / log-cabin lattices of 2 to 4 layers, (f) groove pyramids, cylinders resting in the valley between a parallel pair, and (g) loose criss-cross piles built by dropping cylinders sequentially at random yaw over a small area, matching the reference demo arrangement; shallow-angle multi-contact resting is the dominant contact mode. Regimes (d)--(g) and rolling-to-rest trajectories get heavy weight: cylinders must appear often at rest and mid-roll, since roll-and-settle is the hardest low-energy behavior in this domain and the demo's centerpiece.
  \item Actions: external forces applied at a random surface point of a random object, recorded as the action channel. Two kinds, matching the two interaction modes. \emph{Pokes}: short bursts of random magnitude and direction at 1 to 3 random times. \emph{Grabs}: a spring-damper attached from the surface point to a target that moves along a smooth random path (lift, sideways pull, slide along the table) for 0.5 to 2\,s, then releases. In pile scenes the grabbed object is often load-bearing, which makes the data dense in support-removal collapse, the demo's signature behavior: pull the bottom pencil and the pile falls. The same spring law runs at inference, so grabbing is a force interaction, never a kinematic override, and a pinned pencil physically resists extraction instead of dragging neighbors into interpenetration.
  \item 300 steps recorded at $\Delta t = 1/60$\,s (simulated at 240\,Hz internally, subsampled), storing per-body poses and twists.
\end{itemize}
Target: 5{,}000 trajectories $\rightarrow$ $\sim$1.5M transitions, an hour or two on the M4 Pro's CPU cores with multiprocessing.

\subsection{Particle-ization}
For every trajectory, particles are sampled once on each object's surface mesh with the exact Stage C sampler (farthest-point at fixed spacing), stored as body-frame offsets; world positions at every step derive from the recorded poses. Storing poses plus offsets instead of raw point clouds keeps the dataset in the tens of MB and guarantees rigid-consistent ground truth.

\subsection{Domain randomization for the reconstruction gap}
Reconstructed objects are not perfect primitives. At training time: perturb particle positions with 0.5 to 2\,mm surface noise, randomly delete 5 to 15\% of particles, and jitter per-body counts, so a lumpy reconstructed cluster still looks in-distribution.

%=====================================================================
\section{Training and Evaluation}
\label{sec:training}
%=====================================================================

\begin{center}
\begin{tabular}{@{}ll@{}}
\toprule
Latent width / MP layers & 128 / 10 \\
Parameters & $\approx$ 5M \\
Optimizer & AdamW, lr $10^{-4} \rightarrow 10^{-6}$ (exponential decay) \\
Batch & 2 to 4 scenes' worth of graphs per step \\
Steps & 300k to 500k single-step, then 20k rollout fine-tune \\
Precision & float32 (MPS has no float64) \\
Hardware & M4 Pro MPS; optionally one A100 day ($<\$30$) for the final run \\
Wall clock & 1 to 3 days on MPS \\
\bottomrule
\end{tabular}
\end{center}

Practical notes: run with \texttt{PYTORCH\_ENABLE\_MPS\_FALLBACK=1}; build graphs in dataloader workers, not the training loop; log validation \emph{rollout} error every few thousand steps, since single-step loss is a poor proxy for what matters.

\subsection{Evaluation metrics}
\begin{itemize}[itemsep=2pt]
  \item \textbf{Rollout error:} mean per-body translation and rotation error vs.\ ground truth at 50, 150, 300 steps.
  \item \textbf{Penetration depth:} mean and max interpenetration between bodies and with the ground; the learned model has no hard constraints, so this quantifies how well contact was learned.
  \item \textbf{Stability:} fraction of initially stable scenes, settled criss-cross piles included, that remain at rest after 300 unperturbed steps. A model that knocks over stable piles by itself is not ready.
  \item \textbf{Support-removal fidelity:} scripted extraction of a load-bearing body from a settled pile, run identically through PyBullet and the model; compare which bodies fall and where they settle. This is the demo's signature interaction, so it gets its own number.
  \item \textbf{Energy sanity:} total energy must not grow over a passive rollout.
  \item \textbf{Qualitative:} side-by-side videos, learned vs.\ PyBullet, same seed and same pokes.
\end{itemize}

%=====================================================================
\section{Step-by-Step Build Plan from Zero}
\label{sec:roadmap}
%=====================================================================

Milestones are ordered so that every one ends with something visible and the riskiest work (the learned simulator) happens first; the reconstruction front end, which is mostly running pretrained checkpoints, comes after. Time estimates assume part-time work.

\subsection{Milestone 0: environment and skeleton (half a day)}
\begin{lstlisting}
uv init physsplat && cd physsplat
uv add torch numpy scipy pybullet h5py rerun-sdk hdbscan trimesh onnx onnxruntime
python -c "import torch; print(torch.backends.mps.is_available())"
\end{lstlisting}
\begin{lstlisting}
physsplat/
  datagen/     # PyBullet scene generation, particle sampling
  model/       # GNN, integrators, losses (plain PyTorch)
  train/       # loops, config, checkpoints
  recon/       # image -> splats -> bodies -> particles
  export/      # ONNX export + parity tests
  web/         # Next.js app (Milestone 5)
  scripts/     # entry points
\end{lstlisting}
Install the Rerun viewer early: logging 3D particle states to Rerun is the fastest way to inspect everything this project produces without writing viewer code.

\subsection{Milestone 1: synthetic data you can watch (week 1)}
\begin{enumerate}[itemsep=2pt]
  \item \texttt{datagen/scenes.py}: random scene generator over the primitive-shape family (start with boxes only, add cylinders and convex hulls once the loop works), settle, record poses at 60\,Hz, apply logged impulses. Verify visually in PyBullet's GUI before scaling.
  \item \texttt{datagen/particles.py}: farthest-point fixed-spacing surface sampler over an arbitrary mesh (trimesh handles primitives and hulls uniformly), body-frame offsets. Reused verbatim in Stage C; keep it shared.
  \item HDF5 writer (poses, offsets, actions, material parameters) and a loader that reconstitutes world-space particles.
  \item Generate 50 trajectories, log a few to Rerun, and check: no interpenetration in ground truth, impulses visibly kick objects, scenes settle.
  \item Scale to 5{,}000 trajectories across the full shape family. Compute and store normalization statistics.
\end{enumerate}
\textbf{Deliverable:} a Rerun recording of ground-truth physics from your own generator.

\subsection{Milestone 2: the learned simulator (weeks 2--3, the crux)}
\begin{enumerate}[itemsep=2pt]
  \item Graph builder (grid-hash radius search), unit-tested against brute-force $O(N^2)$.
  \item Encode--Process--Decode with both heads (per-body, and per-particle plus Kabsch). Plain PyTorch.
  \item Overfit a single trajectory to near-zero loss. If this fails, the bug is in features, normalization, or indexing, and no amount of data will save it.
  \item Train on the full set with noise injection. Targets in order: a dropped box lands and stops; a two-box stack stays up; a poked stack topples plausibly; a flicked pencil rolls and comes to rest.
  \item Rollout evaluation harness (metrics of Section~\ref{sec:training}) and side-by-side videos.
  \item Rollout fine-tuning pass.
\end{enumerate}
\textbf{Deliverable:} video, learned vs.\ PyBullet, same poke, visually matching for 5 seconds. Once this exists the project cannot fail completely; everything after is integration.

\subsection{Milestone 3: interactive local demo, synthetic scenes (week 4)}
A minimal Three.js page (objects as plain meshes, no splats yet) connected over WebSocket to a Python process running the model on MPS. Click to poke. This validates the interaction loop, the impulse feature, and reset, with none of the web-inference complexity.
\textbf{Deliverable:} you can knock over a synthetic stack with your mouse, locally.

\subsection{Milestone 4: the reconstruction front end (week 5)}
\begin{enumerate}[itemsep=2pt]
  \item Photograph 3 to 5 distinctly colored objects (blocks and erasers first; add pencils once boxes work), plain background, good light. Run \texttt{rembg}, then TripoSR (easier), LGM if quality disappoints. Export Gaussians or points as \texttt{.ply}.
  \item Stage B: gravity align, RANSAC ground, HDBSCAN with position+color features, per-body OBB, mass, inertia. Tune $\lambda$ and \texttt{min\_cluster\_size} on your test scene. Log colored-by-body clouds to Rerun.
  \item Stage C with the shared sampler from Milestone 1. Verify in Rerun that reconstructed-scene particles look like training particles.
  \item Feed the result to the Milestone 3 demo. Expect one iteration round (scale and gravity mismatches are the usual suspects; domain randomization is the usual fix).
\end{enumerate}
\textbf{Deliverable:} photo in, interactive (mesh-rendered) physics out, locally.

\subsection{Milestone 5: browser inference and splat rendering (weeks 6--7)}
\begin{enumerate}[itemsep=2pt]
  \item ONNX export with symbolic shapes; a parity test asserting PyTorch and onnxruntime agree to $10^{-4}$ over a 100-step rollout, \emph{before} any frontend code touches the model.
  \item Next.js + React Three Fiber app: load the scene packet, ONNX Runtime Web (WebGPU, WASM fallback), TypeScript graph builder and integrator, splat rendering with per-body batches, OBB picking, impulse UI, sliders, reset.
  \item Reconstruction endpoint: FastAPI on Modal, or precompute scene packets and serve them as static files, which removes all server cost and cold starts while reviewers still get full interactivity.
  \item Deploy to Vercel.
\end{enumerate}
\textbf{Deliverable:} public URL. Photo-derived objects, poked in a browser, physics by neural network, client-side.

\subsection{Milestone 6: polish and writeup (week 8)}
Dark-mode UI pass, a 60-second screen recording, a README with the architecture diagram, evaluation table, and honest limitations. Optionally the multi-view capture path for a hero scene.

%=====================================================================
\section{Risks and Fallbacks}
%=====================================================================

\begin{center}
\begin{tabular}{@{}p{5.2cm}p{2.2cm}p{6.6cm}@{}}
\toprule
\textbf{Risk} & \textbf{Likelihood} & \textbf{Fallback} \\
\midrule
Per-body head underfits contact detail & Medium & Per-particle head + Kabsch projection (built in Milestone 2, one config switch) \\
Single-image recon too poor to cluster & Medium & Multi-view phone capture path; or reconstruct each object separately and compose \\
LGM/TripoSR broken on MPS & Medium & Run Stage A on Colab or a rented GPU; it is offline and one-shot \\
ONNX export op gaps & Low--Med & Implement scatter/graph ops in TypeScript, keep only MLPs in ONNX; or WebSocket + local server for live demos and a recorded video for the public page \\
Browser too slow at 60\,Hz & Low & 30\,Hz physics + interpolation; coarser particle spacing \\
Rolling never settles (pencils creep or jitter at rest) & Medium & Overweight low-energy rolling states in the data mix; settle damping below a kinetic-energy threshold; if needed, restrict the public demo to non-rolling shapes \\
Long-rollout drift despite noise injection & Medium & Stronger noise, longer rollout fine-tuning, soft settle damping near zero kinetic energy \\
MPS training too slow & Low & One A100 day, $<\$30$; the code is device-agnostic \\
\bottomrule
\end{tabular}
\end{center}

The compound risk profile is favorable because every medium risk has a fallback that is built into an earlier milestone rather than bolted on later.

%=====================================================================
\section{Budget}
%=====================================================================

\begin{center}
\begin{tabular}{@{}lll@{}}
\toprule
\textbf{Item} & \textbf{Cost} & \textbf{Notes} \\
\midrule
Data generation & \$0 & PyBullet, CPU, local \\
Training & \$0 -- \$25 & MPS free; optional A100 day for the final model \\
Reconstruction & \$0 & Local MPS, or free Colab for LGM \\
Hosting & \$0 & Vercel Hobby; static scene packets; client-side inference \\
Test objects and a phone & already owned & \\
\midrule
\textbf{Total} & \textbf{\$0 -- \$25} & \\
\bottomrule
\end{tabular}
\end{center}

%=====================================================================
\begin{thebibliography}{9}
\bibitem{gns} A.\ Sanchez-Gonzalez, J.\ Godwin, T.\ Pfaff, R.\ Ying, J.\ Leskovec, P.\ Battaglia. \emph{Learning to Simulate Complex Physics with Graph Networks.} ICML 2020.
\bibitem{fignet} K.\ Allen, T.\ Lopez-Guevara, Y.\ Rubanova, et al. \emph{Learning Rigid Dynamics with Face Interaction Graph Networks.} ICLR 2023.
\bibitem{3dgs} B.\ Kerbl, G.\ Kopanas, T.\ Leimk\"uhler, G.\ Drettakis. \emph{3D Gaussian Splatting for Real-Time Radiance Field Rendering.} SIGGRAPH 2023.
\bibitem{lgm} J.\ Tang, Z.\ Chen, X.\ Chen, T.\ Wang, G.\ Zeng, Z.\ Liu. \emph{LGM: Large Multi-View Gaussian Model for High-Resolution 3D Content Creation.} ECCV 2024.
\bibitem{triposr} D.\ Tochilkin, D.\ Pankratz, Z.\ Liu, et al. \emph{TripoSR: Fast 3D Object Reconstruction from a Single Image.} arXiv:2403.02151, 2024.
\end{thebibliography}

\end{document}
